Evaluation of large language model driving agents for intelligent vehicles is constrained not only by task quality, but also by whether local runtime behavior leaves runs comparable in the first place. This paper presents DiLu-Ollama as a local-first evaluation framework for studying fine-tuned LLM driving agents under a shared highway benchmark while treating benchmark validity as part of the measured outcome. The framework combines local inference, task-conditioned highway simulation, and a post-hoc analysis pipeline that records task metrics together with ranking eligibility, invalidity reasons, and latency behavior. Using this framework, we study lightweight, midclass, and highclass model tiers under one runtime regime and find a sharp interpretability boundary: lightweight is the only tier that supports stable comparative ranking, midclass remains screening-only, and all tested 14B models fail ranking eligibility under the current local setup. Within the valid region, the lightweight Phi pair provides the only exact base-versus-fine-tuned comparison and shows that fine-tuning can improve driving score without an obvious mean-latency penalty. These results position benchmark validity as a first-class requirement for automated-vehicle evaluation of local LLM agents and show that fine-tuning claims are only trustworthy where the evaluation regime itself remains stable.
